\chapter{Qwen分层执行与视觉完成检查}

本章沿着一个问题展开：高层到底增加了什么？依次比较翻译、子目标分解、失败恢复、完成判断及视觉干预。不同版本有独立输出与初态，不能用后来的成功改写早期失败。

\section{中文长任务分解与翻译对照}

\label{stage:s09}

让Qwen先理解中文长指令，再把英文子目标交给\pifive{}执行。为区分语言翻译与分层规划的收益，增加只翻译整句指令的平坦执行对照。

\subsection{高层与低层的输入分工}\label{s09-detail-1}

高层Qwen接收单独渲染的512像素第三人称图像和中文任务；低层\(\pi\)0.5仍使用官方外部、腕部图像与机器人状态转换。高层额外渲染不替换或重采样策略输入。Qwen输出英文子目标，\(\pi\)0.5按子目标生成动作块，模型之间没有替换内部网络层。

完成检查比较子目标开始帧与当前帧。物体入篮后可能只露出深色顶部，单帧容易把桌上相似物体误认成目标；前后对比用于关注目的区域变化。夹爪张开门控减少持物状态下的提前确认，连续两次完成答复才切换。门控用的是机器人自身状态，不是物体真值。

最后一个子目标不由Qwen提前终止，继续到官方成功或520步上限，与平坦基线保持相同终止方式。前面子目标的确认只用模型视觉结果；官方关系另存用于分析。后来的同场景视觉停止系统采用另一种停止协议，不能混用两者。

\subsection{中文长任务的版本与强基线}\label{s09-detail-2}

\FloatBarrier\begingroup\small\setlength{\tabcolsep}{2.5pt}
\begin{longtable}[]{@{}
  >{\raggedright\arraybackslash}p{(\linewidth - 6\tabcolsep) * \real{0.2711}}
  >{\raggedleft\arraybackslash}p{(\linewidth - 6\tabcolsep) * \real{0.2119}}
  >{\raggedleft\arraybackslash}p{(\linewidth - 6\tabcolsep) * \real{0.2119}}
  >{\raggedright\arraybackslash}p{(\linewidth - 6\tabcolsep) * \real{0.3050}}@{}}
\caption{中文长任务的版本与强基线：设置与结果}\label{tab:s09-1}\tabularnewline
\toprule\noalign{}
\begin{minipage}[b]{\linewidth}\raggedright
配置
\end{minipage} & \begin{minipage}[b]{\linewidth}\raggedleft
回合数
\end{minipage} & \begin{minipage}[b]{\linewidth}\raggedleft
成功
\end{minipage} & \begin{minipage}[b]{\linewidth}\raggedright
使用范围
\end{minipage} \\
\midrule\noalign{}
\endfirsthead
\toprule\noalign{}
\begin{minipage}[b]{\linewidth}\raggedright
配置
\end{minipage} & \begin{minipage}[b]{\linewidth}\raggedleft
回合数
\end{minipage} & \begin{minipage}[b]{\linewidth}\raggedleft
成功
\end{minipage} & \begin{minipage}[b]{\linewidth}\raggedright
使用范围
\end{minipage} \\
\midrule\noalign{}
\endhead
\bottomrule\noalign{}
\endlastfoot
直接中文给\(\pi\)0.5 & 30 & 22/30 & 平坦中文基线 \\
早期分层v1 & 30 & 16/30 & 无防抖且末目标可提前结束 \\
分层v2 & 30 & 27/30 & 含一次重跑覆盖，作为开发记录 \\
冻结分层v3首次运行 & 30 & 28/30 & 正式保留结果 \\
Qwen只翻译整句后给\(\pi\)0.5 & 30 & 29/30 & 必须保留的强对照 \\
英文原指令平坦执行 & 30 & 29/30 & 早前公开任务基线 \\
\end{longtable}\endgroup

十个任务各三个公开初态。分层比直接中文高，但没有超过仅翻译；这组样本支持语言处理有作用，尚不支持独立的分层成功率优势。早期把收益全部归为语言理解的表述也过强，因为未对所有影响因素作独立消融。

v1到后续版本的改变包括两票确认、最后子目标不提前结束、放宽未使用提示字段的格式检查。多个改动同时发生，因此不是单因素因果实验。某任务曾首跑超时、重跑成功；后续冻结版本保留首次结果，不用重跑覆盖作为正式统计。

\subsection{重规划与整句回退}\label{s09-detail-3}

子目标超预算时，可按当前图像重规划剩余步骤；若新计划原样重复卡住的步骤，则回退整句英文任务执行。旧正式v3默认关闭重规划，恢复试验另用新输出目录。在三个失败条件上的单例救回没有形成完整30回合回归，后续又发现恢复前缀不同，不能与固定起点试验合并。

对外报告应同时列出仅翻译、冻结分层和直接中文，区分``理解指令''\,``推进阶段''\,``修复动作''。本阶段未训练专门的指令分解适配器，Qwen仍为基座。


\subsection{B与8B完成检查诊断}\label{s09-detail-4}

\subsubsection{范围与方法}\label{s09-detail-5}

使用已有 LIBERO-Long 回合的 8 组 512\(\times\)512 前后帧，两模型依次运行同一个完成检查提示词。样例在 8B 加载前固定，包含四个历史误判与四个对照，各四个正、负标签。这是难例诊断，不是随机测试集，也不是任务成功率评测。

两模型均为官方 Instruct 基座，BF16、SDPA、CUDA 单卡、greedy 解码、最多 96 个新 token，无 LoRA、无 CPU 卸载；先单次预热，再保留每例首次输出。8B 下载的四个权重分片已与官方 LFS SHA256 核对。旧引擎和分层执行程序没有改动。

\subsubsection{实测结果}\label{s09-detail-6}

\FloatBarrier\begingroup\small\setlength{\tabcolsep}{2.5pt}
\begin{longtable}[]{@{}
  >{\raggedright\arraybackslash}p{(\linewidth - 4\tabcolsep) * \real{0.4102}}
  >{\raggedleft\arraybackslash}p{(\linewidth - 4\tabcolsep) * \real{0.2949}}
  >{\raggedleft\arraybackslash}p{(\linewidth - 4\tabcolsep) * \real{0.2949}}@{}}
\caption{实测结果：设置与结果}\label{tab:s09-2}\tabularnewline
\toprule\noalign{}
\begin{minipage}[b]{\linewidth}\raggedright
指标
\end{minipage} & \begin{minipage}[b]{\linewidth}\raggedleft
4B
\end{minipage} & \begin{minipage}[b]{\linewidth}\raggedleft
8B
\end{minipage} \\
\midrule\noalign{}
\endfirsthead
\toprule\noalign{}
\begin{minipage}[b]{\linewidth}\raggedright
指标
\end{minipage} & \begin{minipage}[b]{\linewidth}\raggedleft
4B
\end{minipage} & \begin{minipage}[b]{\linewidth}\raggedleft
8B
\end{minipage} \\
\midrule\noalign{}
\endhead
\bottomrule\noalign{}
\endlastfoot
实际参数量 & 4.438 B & 8.767 B \\
与记录谓词一致 & 4/8 & 4/8 \\
合法输出 & 8/8 & 8/8 \\
加载耗时 & 6.15 s & 8.19 s \\
单次检查耗时中位数，含预处理 & 0.527 s & 0.796 s \\
检查阶段 PyTorch 分配峰值 & 8.44 GiB & 16.53 GiB \\
检查阶段 PyTorch 保留峰值 & 8.48 GiB & 16.57 GiB \\
\end{longtable}\endgroup

显存统计不含桌面、CUDA 上下文等其他进程开销，也不是模型加载全程的整机峰值。耗时为当前提示词和短输出的单次实测。

\FloatBarrier\begingroup\small\setlength{\tabcolsep}{2.5pt}
\begin{longtable}[]{@{}
  >{\raggedright\arraybackslash}p{(\linewidth - 6\tabcolsep) * \real{0.3014}}
  >{\raggedright\arraybackslash}p{(\linewidth - 6\tabcolsep) * \real{0.2329}}
  >{\raggedright\arraybackslash}p{(\linewidth - 6\tabcolsep) * \real{0.2329}}
  >{\raggedright\arraybackslash}p{(\linewidth - 6\tabcolsep) * \real{0.2329}}@{}}
\caption{实测结果：设置与结果}\label{tab:s09-3}\tabularnewline
\toprule\noalign{}
\begin{minipage}[b]{\linewidth}\raggedright
样例
\end{minipage} & \begin{minipage}[b]{\linewidth}\raggedright
记录标签
\end{minipage} & \begin{minipage}[b]{\linewidth}\raggedright
4B
\end{minipage} & \begin{minipage}[b]{\linewidth}\raggedright
8B
\end{minipage} \\
\midrule\noalign{}
\endfirsthead
\toprule\noalign{}
\begin{minipage}[b]{\linewidth}\raggedright
样例
\end{minipage} & \begin{minipage}[b]{\linewidth}\raggedright
记录标签
\end{minipage} & \begin{minipage}[b]{\linewidth}\raggedright
4B
\end{minipage} & \begin{minipage}[b]{\linewidth}\raggedright
8B
\end{minipage} \\
\midrule\noalign{}
\endhead
\bottomrule\noalign{}
\endlastfoot
wrong\_plate & 否 & 否 & 否 \\
false\_done\_mug & 否 & 是 & 否 \\
plate\_viewpoint & 是 & 否 & 是 \\
basket\_occlusion & 是 & 否 & 否 \\
soup\_done & 是 & 是 & 是 \\
soup\_not\_done & 否 & 否 & 是 \\
false\_done\_pot & 否 & 是 & 是 \\
microwave\_done & 是 & 是 & 否 \\
\end{longtable}\endgroup

\subsubsection{解释与限制}\label{s09-detail-7}

该次 8B 修正了 false\_done\_mug 和 plate\_viewpoint 两例，但在 soup\_not\_done 和 microwave\_done 两例上退化，总一致数与 4B 相同；短输出检查耗时中位数增加约 51\%。这组结果没有支持直接替换现有 4B，先保留原在线配置。更大参数量不保证这类完成判断更可靠。

8B BF16 在本机 RTX 4090 上已完成独立加载和双图推理。这次没有让它与 \(\pi\)0.5 同驻留；4B 与 \(\pi\)0.5 的既有显存记录不能直接替代 8B 共驻测试。若要替换在线系统，仍需另测量化或分时运行。

标签来自旧回合中对应首子目标的官方谓词，不放入模型提示词。左右参照、物体遮挡、是否仍被夹爪支撑、落在盘中的接触边界等，在单一视角下可能模糊；尤其 plate\_viewpoint 与 wrong\_plate 两例外观接近而标签不同。应结合原图与完整轨迹解释，不能把每次不一致都当作纯视觉识别错误。

未运行新仿真回合、重规划回归、微调或实体机械臂。8 个诊断样例不足以证明更大的模型能提高最终任务成功率；下一步应先审阅左右参照与遮挡样例，再决定是否做完整配对回归。


\stagefigure{figures/ch04/s09_long0_initial.png}{LIBERO长任务的观察画面，作为高层分解与检查的输入背景。}{fig:s09-retained}

\paragraph{阶段结论}冻结分层版本为28/30，直接中文22/30，仅翻译29/30。结果表明中文收益主要可由语言处理解释，尚未证明分层具有额外成功率优势。 每任务仅三个初态，结论受样本规模限制。历史开发版本和重跑结果应与冻结版本分开。

\stagecode{s09}

\stagevideo{videos/ch04/hierarchical_long_example.mp4}

\FloatBarrier

\section{执行状态管理与受控恢复}

\label{stage:s10}

分析近终点失败、步骤超时和重规划覆盖不足，再通过固定恢复起点、受控扰动和不同恢复方式进行对照。

\subsection{已有结果与该阶段问题}\label{s10-detail-1}

此前已在同一LIBERO/Panda环境中完成真正的VLA闭环：RGB、语言及机器人状态进入官方预训练\(\pi\)0.5检查点，模型输出动作，由模拟器中的控制器执行；不是预录轨迹或直接设置末端位姿。四任务集基线为117/120，另有126次输入和执行设置的配对干预。这些已有结果不计作该阶段新增实验。

中文Long任务的现有结果如下。每组为10任务、每任务3个公开初态，样本量较小。

\FloatBarrier\begingroup\small\setlength{\tabcolsep}{2.5pt}
\begin{longtable}[]{@{}
  >{\raggedright\arraybackslash}p{(\linewidth - 4\tabcolsep) * \real{0.3356}}
  >{\raggedleft\arraybackslash}p{(\linewidth - 4\tabcolsep) * \real{0.2676}}
  >{\raggedright\arraybackslash}p{(\linewidth - 4\tabcolsep) * \real{0.3968}}@{}}
\caption{已有结果与该阶段问题：设置与结果}\label{tab:s10-1}\tabularnewline
\toprule\noalign{}
\begin{minipage}[b]{\linewidth}\raggedright
配置
\end{minipage} & \begin{minipage}[b]{\linewidth}\raggedleft
成功数
\end{minipage} & \begin{minipage}[b]{\linewidth}\raggedright
使用说明
\end{minipage} \\
\midrule\noalign{}
\endfirsthead
\toprule\noalign{}
\begin{minipage}[b]{\linewidth}\raggedright
配置
\end{minipage} & \begin{minipage}[b]{\linewidth}\raggedleft
成功数
\end{minipage} & \begin{minipage}[b]{\linewidth}\raggedright
使用说明
\end{minipage} \\
\midrule\noalign{}
\endhead
\bottomrule\noalign{}
\endlastfoot
\(\pi\)0.5直接接收中文 & 22/30 & 平坦中文基线 \\
Qwen分层v1 & 16/30 & 历史开发版，含2次程序错误 \\
Qwen分层v2 & 27/30 & 含重跑，不作为冻结首次运行结果 \\
Qwen分层v3 & 28/30 & 冻结首次运行结果 \\
Qwen仅翻译，再由\(\pi\)0.5执行 & 29/30 & 必须保留的强对照 \\
\(\pi\)0.5英文原指令 & 29/30 & 来自此前英文基线 \\
\end{longtable}\endgroup

中文直接执行与仅翻译的30对回合中，22对均成功、7对由失败变成功、1对均失败。仅翻译与分层v3之间，28对均成功、1对仅翻译成功、1对均失败。当前结果尚未显示分层超过仅翻译；不能将分层相对中文直接执行的变化全部归因于规划能力，也不能从这组小样本断言分层在其他条件下无价值。

历史审计发现，v2有3次因预算耗尽切换子目标，v3有2次；v3的两次均发生在最终失败回合。超时切换是已确认的程序行为，但不能仅凭相关性断言修掉它就能救回任务。

相关代码在 \path|recovery_study/scripts/audit.py|；逐回合索引在 \path|episode_index.csv|，汇总在 \texttt{audit.json}。

\subsection{执行逻辑改进}\label{s10-detail-2}

原来的执行链为``分解任务---执行子目标---Qwen比较前后帧---切换下一步''。该阶段保留历史入口，新增独立改进版，主要改动如下。

\FloatBarrier\begingroup\small\setlength{\tabcolsep}{2.5pt}
\begin{longtable}[]{@{}
  >{\raggedright\arraybackslash}p{(\linewidth - 2\tabcolsep) * \real{0.4176}}
  >{\raggedright\arraybackslash}p{(\linewidth - 2\tabcolsep) * \real{0.5824}}@{}}
\caption{执行逻辑改进：设置与结果}\label{tab:s10-2}\tabularnewline
\toprule\noalign{}
\begin{minipage}[b]{\linewidth}\raggedright
问题
\end{minipage} & \begin{minipage}[b]{\linewidth}\raggedright
改进
\end{minipage} \\
\midrule\noalign{}
\endfirsthead
\toprule\noalign{}
\begin{minipage}[b]{\linewidth}\raggedright
问题
\end{minipage} & \begin{minipage}[b]{\linewidth}\raggedright
改进
\end{minipage} \\
\midrule\noalign{}
\endhead
\bottomrule\noalign{}
\endlastfoot
子目标超时可能直接跳到下一步 & 超时只记录未完成并触发恢复；没有恢复预算时继续当前目标，直到总步数上限 \\
重规划把所有前序步骤当作完成 & 仅传入被视觉确认的步骤，当前尝试过的目标不加入完成列表 \\
同一边界同时满足完成和超时 & 先处理确认完成，不再对已完成目标重规划 \\
``拿起''与``放下''拆开，但检查要求夹爪张开 & 要求搬运步骤包含放置终态；拒绝常见的独立拿起、抓住、抬起表达 \\
重规划遗漏尚未完成的要求 & 为未确认目标分配编号，返回计划必须覆盖每个编号一次 \\
旋钮操作也要求夹爪张开才检查 & 增加按动作类型检查的实验开关；补测出现误确认，默认不启用 \\
新实验覆盖已有输出 & 新入口拒绝覆盖已有 \texttt{result.json}；原结果与旧入口不修改 \\
\end{longtable}\endgroup

目标编号校验保证的是形式覆盖，不保证自然语言改写的语义一定正确；独立拿起检查也只是针对常见表达的接口保护，不是完整的自然语言证明。视觉确认同样可能错误，因此``已确认''并不等于模拟器真值。

新版本没有将整句重试隐藏在重规划内部。整句重试和重规划作为独立策略比较，避免把重新执行完整任务的收益记到规划器名下。最后一个子目标仍使用LIBERO官方成功或总步数上限结束，因此当前仍是基准评测方案，不代表已解决真机自主终止。

相关代码在 \path|scripts/subgoal_state.py|、\path|scripts/run_hierarchical_recovery.py|、\path|scripts/qwen_planner_recovery.py|。原 \path|run_hierarchical.py|、\path|qwen_planner.py|、\texttt{evaluate.py} 的哈希已核对，未修改。

\subsection{配对恢复实验}\label{s10-detail-3}

\subsubsection{实验起点}\label{s10-detail-4}

固定选择此前v2的三个失败回合：t2初态2（开炉子并放壶）、t4初态2（两杯分别放盘）、t8初态2（两壶放炉台）。每次从同一个公开初态开始，重放该历史回合前150步动作，再进行恢复。每个状态使用31001、31002两组动作噪声。

每次恢复前核对完整模拟器状态哈希。同一任务的各策略和噪声重复均得到相同起点；重放第150步末端位置与旧记录的最大差值为0。每种方式再给370步动作预算，总计不超过原任务520步。策略每次预测10步，只执行前5步后重新观察，配对分支使用相同噪声种子。

这构成三个独立选定故障状态、每状态两次随机重复，不是六个独立场景，更不是新的Long总体成功率评测。不同策略的阶段切换和调用次数可能不同，不能将同种子理解为整个后续轨迹相同。

\subsubsection{先导对照及后续约束版}\label{s10-detail-5}

第一阶段先固定三种策略、共18次尝试：继续原子目标、直接执行官方完整英文任务、原Qwen重规划。原重规划在t8的两次输出结构不合法，这两次未执行恢复动作，保留为程序失败；查明原因后继续剩余预定项，没有重跑覆盖。

随后针对已经发现的目标遗漏与步骤粒度问题，增加约束版重规划，共6次。这是看过先导结果后的探索性改进，不是预先独立测试集。

\FloatBarrier\begingroup\small\setlength{\tabcolsep}{2.5pt}
\begin{longtable}[]{@{}
  >{\raggedright\arraybackslash}p{(\linewidth - 8\tabcolsep) * \real{0.2980}}
  >{\raggedleft\arraybackslash}p{(\linewidth - 8\tabcolsep) * \real{0.1755}}
  >{\raggedleft\arraybackslash}p{(\linewidth - 8\tabcolsep) * \real{0.1755}}
  >{\raggedleft\arraybackslash}p{(\linewidth - 8\tabcolsep) * \real{0.1755}}
  >{\raggedleft\arraybackslash}p{(\linewidth - 8\tabcolsep) * \real{0.1755}}@{}}
\caption{先导对照及后续约束版：设置与结果}\label{tab:s10-3}\tabularnewline
\toprule\noalign{}
\begin{minipage}[b]{\linewidth}\raggedright
恢复方式
\end{minipage} & \begin{minipage}[b]{\linewidth}\raggedleft
尝试数
\end{minipage} & \begin{minipage}[b]{\linewidth}\raggedleft
实际执行动作
\end{minipage} & \begin{minipage}[b]{\linewidth}\raggedleft
成功
\end{minipage} & \begin{minipage}[b]{\linewidth}\raggedleft
程序失败
\end{minipage} \\
\midrule\noalign{}
\endfirsthead
\toprule\noalign{}
\begin{minipage}[b]{\linewidth}\raggedright
恢复方式
\end{minipage} & \begin{minipage}[b]{\linewidth}\raggedleft
尝试数
\end{minipage} & \begin{minipage}[b]{\linewidth}\raggedleft
实际执行动作
\end{minipage} & \begin{minipage}[b]{\linewidth}\raggedleft
成功
\end{minipage} & \begin{minipage}[b]{\linewidth}\raggedleft
程序失败
\end{minipage} \\
\midrule\noalign{}
\endhead
\bottomrule\noalign{}
\endlastfoot
继续原子目标，确认后切换 & 6 & 6 & 0 & 0 \\
整句英文任务重试 & 6 & 6 & 0 & 0 \\
原Qwen重规划，无整句回退 & 6 & 4 & 0 & 2 \\
带目标覆盖和终态约束的重规划 & 6 & 6 & 0 & 0 \\
炉台按动作类型放开检查门控（追加探测） & 2 & 2 & 0 & 0 \\
\end{longtable}\endgroup

约束版6次均返回可接受的终态计划，保留未确认目标，格式错误不再出现；但没有救回任务。接口正确性改善与任务成功率改善必须分开描述。

\stagefigure{figures/ch04/s10_recovery_outcomes_dea9060e.png}{三个恢复起点及不同方式的结束画面}{fig:s10-detail-1}

图中每行为同一个恢复起点；后续列为继续、整句重试和约束重规划的首组噪声结果。图片展示结束状态，分数仍由官方判据计算，不按人工观感判成功。

\subsubsection{正常初态和指定纠偏指令}\label{s10-detail-6}

新执行入口另做两个正常初态功能检查：t0初态0完成，292步成功；t9初态0在520步超时，最终杯子入炉判据为假、门关闭判据为真，合计1/2。该配置同时改变了规划提示词、超时规则和预算，不能用这两例隔离某一改动的影响，也不能据此宣布新版优于旧版。

针对白杯偏心问题，另做两次人工指定指令探测，明确要求``重新拿起白杯，放到左盘中心，再将黄白杯放到右盘中心''。不向策略提供目标坐标或数值阈值，两次均未完成。它只能说明这两个噪声重复下，明确语言要求仍未解决该状态的纠偏；不是Qwen自主生成指令的结果，也不足以证明\(\pi\)0.5普遍无法纠偏。

该阶段共26次恢复尝试、2次指定指令探测、2次正常初态检查；其中2次因规划格式错误未执行动作。实际新分支和正常初态共推进10432步，另有前缀重放及接触诊断重放，不将它们重复计作新模型回合。

相关代码在 \path|recovery_study/scripts/run_paired.py|、\path|run_contract.py|、\path|run_guided.py|。四份协议文件分别保留各阶段实验设计；门控补测见 \path|typed_gate_protocol.json| 与 \path|run_typed_gate.py|。

\subsection{具体失败分析}\label{s10-detail-7}

\subsubsection{白杯：已接触目标盘，但不够居中}\label{s10-detail-8}

此前将该状态描述为``放错盘''不准确。固定重放并检查目标盘的接触与位置关系后发现：白杯接触的是目标盘，杯心也高于盘心，但水平中心距离约38.5毫米。LIBERO的该 \texttt{On} 判据还要求中心距离小于30毫米，因此失败原因是偏心，而不是盘子身份错误。

\stagefigure{figures/ch04/s10_plate_boundary.png}{白杯接触与判定边界}{fig:s10-detail-2}

继续、整句重试和原重规划的首组结果中，白杯偏差均保持在约38.5毫米；黄白杯则已进入其目标盘的判定范围。即使任务在语义上``看起来完成了''，剩余的几何误差仍需要策略主动纠正。仅让高层重复原要求并没有产生这次纠偏。

相关代码在 \path|recovery_study/scripts/inspect_contact.py|，数值在 \path|plate_contact_diagnostic.json|。这些位置、接触和谓词只用于离线解释，不输入控制策略或Qwen。

\subsubsection{炉台：重规划可能省略尚未完成的要求}\label{s10-detail-9}

原重规划在两个重复中都只给出放壶步骤，没有保留未确认的开炉子要求。约束版保留了开炉子和放壶两个目标，但任务仍超时。

继续原目标和约束版的一些轨迹中，开炉子判据曾经成立；系统却没有询问完成检查。原因是机器人仍夹着旋钮，而统一的搬运松爪门控屏蔽了查询。例如约束版首组在恢复后第165步已达到开炉子判据，但该回合零次完成查询。这说明不能把所有失败都归为执行策略不会开炉子。

据此追加两次配对探测：状态、动作噪声、Qwen生成的计划和提示词均与约束版一致，仅允许开关炉子等状态变化目标在夹爪闭合时查询，搬运目标仍要求释放。两次均失败，且分别在恢复后第80步、第140步被两票误判完成；当时开炉子真值仍为假，流程已经切换。门控过严可能漏掉有效完成，放开门控又暴露检查器误判，不能靠简单开关解决。

该功能保留为 \texttt{-\/-operation-checks} 实验开关，默认关闭。它验证的是门控与视觉判定之间的耦合问题，不作为成功率改进。真值只用于离线诊断和记录，运行切换仍由视觉回答决定。

\subsubsection{双壶：格式合法还需要步骤可检查}\label{s10-detail-10}

原重规划两次返回字符串列表，旧接口要求对象列表，导致程序失败；即使将其机械转换成对象，内容仍是独立``拿起---放下''步骤，与张开夹爪后检查搬运终态的规则不匹配。约束版不再产生这些接口错误，但仍未把两壶放好。接受JSON不是完成任务，放宽解析也不能代替正确的技能粒度。

\subsubsection{微波炉：两票仍可能重复同一个误判}\label{s10-detail-11}

正常初态检查中，第300步与第320步的视觉回答均认为杯子已放入微波炉；这两个时刻对应的官方入炉判据均为假。第二票导致切换到关门阶段，最终门关闭但杯子目标未完成。

连续两票能过滤部分瞬时判断错误，却不是两份独立信息。模型在相近视角下可能重复同一个判断错误。此前对v2/v3的t2、t4首阶段共52次历史查询做了人工对齐的诊断统计，虽未见确认后错误切换，也不能推出其他任务安全；该阶段微波炉回合就是明确的反例。

\subsubsection{为什么旧恢复成功例与这次结果不同}\label{s10-detail-12}

此前重规划先导曾在t4初态2成功，但不能视为该阶段同一恢复起点的重复。对比旧先导与该次固定来源回合，第150步末端位置相差约37.1毫米，前150步动作也不相同。该阶段固定失败状态的比较回答的是条件恢复能力，不能与不同轨迹起点直接合并计算。

\subsection{验证、限制与后续工作}\label{s10-detail-13}

已通过10项状态与门控测试和6项目标约束测试；26次恢复尝试的同任务起点哈希一致，末端前缀重放误差为0。原正式结果和旧运行入口保持不变。所有新动作都在仿真中执行，没有连接实体机械臂；该阶段没有新权重下载、训练或量化。

目前最值得继续的方向不是扩大模型，而是两个较小、可明确验证的问题。

\textbf{首先，提高完成检查的可靠性。} 从现有轨迹中整理``已接触但偏心、靠近容器但未放入、仍被夹持、旋钮尚未开启、目标遮挡''等样例，按完整回合划分开发和独立检查集。目标不是只输出成功/失败，而是允许无法判断，并区分目标身份、位置关系和是否已释放。需要同时考察误确认、漏确认和切换延迟；仅提高逐帧正确率不一定提高最终成功率。双视角可作为候选，但尚未验证。

\textbf{其次，单独验证近终点纠偏技能。} 将``从正常起点完成任务''和``已经放偏后重新拿起、居中、释放''分开。先在仿真中建立少量可重复的偏置状态和对应纠偏示教，再判断是否需要针对该技能微调策略。当前两次指定指令失败说明值得补这部分数据，但并非足以启动大规模训练的依据。

恢复能力只有在这些受控状态上稳定出现，再回到10任务\(\times\)3初态的完整回归，检查正常任务是否退化。现在不开展200回合扩大评测，不把1/2的功能检查当作新版通过，也不将新增接口保护包装成成功率提升。

导航、Jetson和SO101迁移继续列为独立验证范围；该阶段没有将它们并入已经完成的LIBERO闭环。


\subsection{后续的受控故障注入与退化复查}\label{s10-detail-14}

\subsubsection{t9定位与单入口}\label{s10-detail-15}

固定t9初态0、相同动作噪声、512像素监测图、双票、夹爪门控及520步上限。初始规划使用已保存的英文子目标，消除初次Qwen生成文本变化，再分别替换执行逻辑、物体称呼和恢复设置。

\FloatBarrier\begingroup\small\setlength{\tabcolsep}{2.5pt}
\begin{longtable}[]{@{}
  >{\raggedright\arraybackslash}p{(\linewidth - 4\tabcolsep) * \real{0.4800}}
  >{\raggedright\arraybackslash}p{(\linewidth - 4\tabcolsep) * \real{0.2600}}
  >{\raggedleft\arraybackslash}p{(\linewidth - 4\tabcolsep) * \real{0.2600}}@{}}
\caption{t9定位与单入口：设置与结果}\label{tab:s10-4}\tabularnewline
\toprule\noalign{}
\begin{minipage}[b]{\linewidth}\raggedright
配置
\end{minipage} & \begin{minipage}[b]{\linewidth}\raggedright
结果
\end{minipage} & \begin{minipage}[b]{\linewidth}\raggedleft
步数
\end{minipage} \\
\midrule\noalign{}
\endfirsthead
\toprule\noalign{}
\begin{minipage}[b]{\linewidth}\raggedright
配置
\end{minipage} & \begin{minipage}[b]{\linewidth}\raggedright
结果
\end{minipage} & \begin{minipage}[b]{\linewidth}\raggedleft
步数
\end{minipage} \\
\midrule\noalign{}
\endhead
\bottomrule\noalign{}
\endlastfoot
历史逻辑＋原``yellow and white mug''计划 & 成功 & 257 \\
严格逻辑＋同一原计划 & 成功 & 257 \\
严格逻辑＋新``mug with yellow handle''计划 & 成功 & 267 \\
严格逻辑＋新计划＋150步触发恢复 & 成功 & 267 \\
\end{longtable}\endgroup

前两组在相同原计划下的动作轨迹逐点一致。这次未稳定复现之前的超时，不能指定某个改动就是原因。旧失败回合与该次相同配置的初始图像逐像素一致，但第一步动作最大绝对差约0.000247，后续轨迹已经分开。四组结果及对比保存在\path|t9_results.json|、\path|t9_conclusion.json|；记录为未定位，不继续展开排查。

对外执行入口统一为\path|scripts/run_hierarchical.py|：

\begin{itemize}
\tightlist
\item
  \path!--execution-profile legacy|strict!选择历史或严格状态转换。
\item
  \path!--planner-profile base|contract!选择原始或目标覆盖/终态约束规划。
\item
  \path!--recovery-mode continue|full|replan!区分继续、整句重试和重规划。
\item
  \texttt{-\/-operation-checks}仍是默认关闭的实验选项。
\end{itemize}

Qwen worker统一由\path!qwen_planner.py --profile base|contract!启动。原recovery文件保留实现模块和兼容包装，共享图像、串口无关的进程通信等帮助函数，不再维护两套worker循环。历史算法实现保留用于复现；后续不增加平行启动入口。纯平坦执行也不再加载无用的Qwen worker。

\subsubsection{受控扰动实验}\label{s10-detail-16}

\paragraph{三类扰动}\label{s10-detail-17}

\FloatBarrier\begingroup\small\setlength{\tabcolsep}{2.5pt}
\begin{longtable}[]{@{}
  >{\raggedright\arraybackslash}p{(\linewidth - 4\tabcolsep) * \real{0.2717}}
  >{\raggedright\arraybackslash}p{(\linewidth - 4\tabcolsep) * \real{0.3145}}
  >{\raggedright\arraybackslash}p{(\linewidth - 4\tabcolsep) * \real{0.4138}}@{}}
\caption{三类扰动：设置与结果}\label{tab:s10-5}\tabularnewline
\toprule\noalign{}
\begin{minipage}[b]{\linewidth}\raggedright
类型
\end{minipage} & \begin{minipage}[b]{\linewidth}\raggedright
Long任务
\end{minipage} & \begin{minipage}[b]{\linewidth}\raggedright
注入方式
\end{minipage} \\
\midrule\noalign{}
\endfirsthead
\toprule\noalign{}
\begin{minipage}[b]{\linewidth}\raggedright
类型
\end{minipage} & \begin{minipage}[b]{\linewidth}\raggedright
Long任务
\end{minipage} & \begin{minipage}[b]{\linewidth}\raggedright
注入方式
\end{minipage} \\
\midrule\noalign{}
\endhead
\bottomrule\noalign{}
\endlastfoot
撤销一件入篮物体 & t0，汤罐和番茄酱入篮 & 第一个目标放置后，将汤罐恢复到原桌面姿态 \\
破坏操作前提 & t3，碗入底层抽屉并关上 & 在放碗前关闭初始打开的底层抽屉 \\
撤销一个杯子的放置 & t4，双杯分别放盘 & 白杯放置并释放后，将它恢复到原桌面姿态 \\
\end{longtable}\endgroup

每类使用公开初态0/1/2。对同一注入状态，比较继续原计划、整句英文重试、Qwen基于当前图像重规划。三分支的扰动状态哈希完全一致，动作噪声种子相同，各给420步恢复预算。t4部分源计划只有一个整句目标，此时``继续''保留真实的当前整句指令，没有为了制造差距而改成故意不完整的指令。

环境代码使用真值构造和检查扰动；策略和Qwen不读取这些值。恢复调用发生在约定的注入时刻，测试的是已通知扰动后的恢复能力，不是自动故障检测。新分支的恢复预算也不作为标准Long排行榜分数。

\paragraph{结果}\label{s10-detail-18}

\FloatBarrier\begingroup\small\setlength{\tabcolsep}{2.5pt}
\begin{longtable}[]{@{}
  >{\raggedright\arraybackslash}p{(\linewidth - 6\tabcolsep) * \real{0.2603}}
  >{\raggedleft\arraybackslash}p{(\linewidth - 6\tabcolsep) * \real{0.2466}}
  >{\raggedleft\arraybackslash}p{(\linewidth - 6\tabcolsep) * \real{0.2466}}
  >{\raggedleft\arraybackslash}p{(\linewidth - 6\tabcolsep) * \real{0.2466}}@{}}
\caption{结果：设置与结果}\label{tab:s10-6}\tabularnewline
\toprule\noalign{}
\begin{minipage}[b]{\linewidth}\raggedright
扰动
\end{minipage} & \begin{minipage}[b]{\linewidth}\raggedleft
继续原计划
\end{minipage} & \begin{minipage}[b]{\linewidth}\raggedleft
整句重试
\end{minipage} & \begin{minipage}[b]{\linewidth}\raggedleft
Qwen重规划
\end{minipage} \\
\midrule\noalign{}
\endfirsthead
\toprule\noalign{}
\begin{minipage}[b]{\linewidth}\raggedright
扰动
\end{minipage} & \begin{minipage}[b]{\linewidth}\raggedleft
继续原计划
\end{minipage} & \begin{minipage}[b]{\linewidth}\raggedleft
整句重试
\end{minipage} & \begin{minipage}[b]{\linewidth}\raggedleft
Qwen重规划
\end{minipage} \\
\midrule\noalign{}
\endhead
\bottomrule\noalign{}
\endlastfoot
t0，撤销入篮 & 0/3 & 2/3 & 0/3 \\
t3，关闭抽屉 & 0/3 & 0/3 & 0/3 \\
t4，撤销白杯 & 2/3 & 1/3 & 0/3 \\
合计 & 2/9 & 3/9 & 0/9 \\
\end{longtable}\endgroup

\stagefigure{figures/ch04/s10_recovery_pair_ae5ee43c.png}{同一扰动状态下的恢复结果}{fig:s10-detail-3}

有控制策略救回的例子，说明该阶段不全是不可执行状态；但抽屉类仍未被任一方法救回，不能把全部九个状态预设为可恢复。Qwen方案未超过整句重试，因此该阶段恢复实验收尾，不继续扩大初态数。

保留的规划文本也暴露了两个具体问题：入篮任务可能添加未要求的开罐/开瓶动作；抽屉关闭后可能仍只给出放碗和关抽屉，没有补打开这个前提。它们不是JSON格式检查能解决的语义问题，该阶段不再追加规划器分支。

协议在\path|injection_protocol.json|，逐例结果在\path|injection_results.json|，视频、状态和动作在\path|results/injection/|。


\paragraph{阶段结论}三个自然失败状态的26次恢复尝试未成功；后续27次受控恢复也未证实Qwen带来额外收益。分析修正了对物体偏心与视觉完成状态的早期解释。 恢复起点选择会影响实验难度。未成功的结果同样保留，但不把诊断本身写成恢复性能提升。

\stagecode{s10}

\FloatBarrier

\section{视觉完成检查器的低秩适配}

\label{stage:s11}

把问题转为开始画面、当前画面和子目标指令到完成判断。构造偏心、悬空、夹持未放等近似完成样例，训练并验证检查器。

\subsection{首版检查器：数据与阈值校准}\label{s11-detail-1}

\subsubsection{完成检查器LoRA}\label{s11-detail-2}

\paragraph{数据与标签}\label{s11-detail-3}

从此前四个LIBERO任务集的120个基线回合提取720组图像对。第一张为回合初始图，第二张取早期、中段、终点前20步、终点前5步、终点前1步和终点。视频第k帧对应执行k步后的状态，标签读取\texttt{trace.predicate{[}k-1{]}}，不是Qwen生成答案。

按初态编号隔离：初态0训练、初态1验证、初态2测试，各40个完整回合、240例。回合和初态哈希没有交叉，完全相同的图像对及指令没有跨集合重叠。任务本身在三个集合中共享，这验证的是该次LoRA的初态留出效果，不是新任务泛化；公共初态也不等于基座预训练从未见过。

终点前5步和前1步共240例被单列为边界样例。其中训练侧80例不参加训练，验证侧80例不参加权重选择；测试侧80例仍完整报告，不隐藏难例。因此实际训练160例，权重选择使用验证集中的固定32正、32负共64例，最终测试240例分为常规160例和边界80例。后续阈值校准使用完整验证集240例，测试分数不用于选择阈值。

图像统一为两张256\(\times\)256 RGB，与源视频分辨率一致；该阶段没有深度输入。标签表示完整任务/目标条件全部满足。普通任务提供位置关系样例，Long任务提供多条件样例；未把最终判据随意映射给任意中间子目标。

\paragraph{训练配置与权重}\label{s11-detail-4}

基座为固定revision的\path|Qwen/Qwen3-VL-4B-Instruct|，BF16；视觉模块冻结，仅在语言层q/v投影上训练LoRA，rank8、alpha16、dropout0.05。可训练参数2,949,120，AdamW学习率5e-5，梯度累积2，正负均衡抽样，固定180次更新。

每60次更新在同一验证子集评估，以平衡准确率最高、再以误报率更低选择。第180步权重被选中，验证平衡准确率87.5\%。测试集不选权重，没有做多轮超参数搜索。

权重在\path|results/checker_lora/best_adapter/|。训练、验证及微调后测试计时约103秒，不含之前的基座测试、数据提取、初次加载和后续校准。PyTorch记录的峰值分配约8.79GiB。独立新进程已重新载入保存的adapter，独立CLI输出与保存的测试概率一致。

\paragraph{独立测试：常规识别提升，边界误报严重}\label{s11-detail-5}

\FloatBarrier\begingroup\small\setlength{\tabcolsep}{2.5pt}
\begin{longtable}[]{@{}
  >{\raggedright\arraybackslash}p{(\linewidth - 4\tabcolsep) * \real{0.3930}}
  >{\raggedleft\arraybackslash}p{(\linewidth - 4\tabcolsep) * \real{0.2965}}
  >{\raggedleft\arraybackslash}p{(\linewidth - 4\tabcolsep) * \real{0.3105}}@{}}
\caption{独立测试：常规识别提升，边界误报严重：设置与结果}\label{tab:s11-1}\tabularnewline
\toprule\noalign{}
\begin{minipage}[b]{\linewidth}\raggedright
240例完整测试，贪心输出
\end{minipage} & \begin{minipage}[b]{\linewidth}\raggedleft
基座
\end{minipage} & \begin{minipage}[b]{\linewidth}\raggedleft
LoRA
\end{minipage} \\
\midrule\noalign{}
\endfirsthead
\toprule\noalign{}
\begin{minipage}[b]{\linewidth}\raggedright
240例完整测试，贪心输出
\end{minipage} & \begin{minipage}[b]{\linewidth}\raggedleft
基座
\end{minipage} & \begin{minipage}[b]{\linewidth}\raggedleft
LoRA
\end{minipage} \\
\midrule\noalign{}
\endhead
\bottomrule\noalign{}
\endlastfoot
总准确率 & 82.9\% & 63.3\% \\
平衡准确率 & 51.6\% & 72.9\% \\
已完成召回，39例 & 2/39（5.1\%） & 34/39（87.2\%） \\
未完成误报，201例 & 4/201（2.0\%） & 83/201（41.3\%） \\
输出格式合法 & 240/240 & 240/240 \\
\end{longtable}\endgroup

基座大多回答未完成，因负例较多而得到较高总准确率。LoRA能识别更多完成样本，但相应误报增加，不能只展示平衡准确率提升。

\FloatBarrier\begingroup\small\setlength{\tabcolsep}{2.5pt}
\begin{longtable}[]{@{}
  >{\raggedright\arraybackslash}p{(\linewidth - 4\tabcolsep) * \real{0.3940}}
  >{\raggedleft\arraybackslash}p{(\linewidth - 4\tabcolsep) * \real{0.2842}}
  >{\raggedleft\arraybackslash}p{(\linewidth - 4\tabcolsep) * \real{0.3217}}@{}}
\caption{独立测试：常规识别提升，边界误报严重：设置与结果}\label{tab:s11-2}\tabularnewline
\toprule\noalign{}
\begin{minipage}[b]{\linewidth}\raggedright
分层查看，贪心输出
\end{minipage} & \begin{minipage}[b]{\linewidth}\raggedleft
基座
\end{minipage} & \begin{minipage}[b]{\linewidth}\raggedleft
LoRA
\end{minipage} \\
\midrule\noalign{}
\endfirsthead
\toprule\noalign{}
\begin{minipage}[b]{\linewidth}\raggedright
分层查看，贪心输出
\end{minipage} & \begin{minipage}[b]{\linewidth}\raggedleft
基座
\end{minipage} & \begin{minipage}[b]{\linewidth}\raggedleft
LoRA
\end{minipage} \\
\midrule\noalign{}
\endhead
\bottomrule\noalign{}
\endlastfoot
常规160例准确率 & 76.9\% & 86.3\% \\
常规负例误报率 & 0/121 & 17/121（14.0\%） \\
边界80例误报率 & 4/80（5.0\%） & 66/80（82.5\%） \\
\end{longtable}\endgroup

\stagefigure{figures/ch04/s11_checker_core_3bdc7ffb.png}{常规160例的混淆矩阵}{fig:s11-detail-1}

\paragraph{验证集阈值校准}\label{s11-detail-6}

另外按固定JSON前缀后true/false的相对概率做受限二分类。基座与LoRA使用同一方法，在验证集误报不超过5\%的约束下选择召回最高的阈值。此项是追加的运行点校准，原始贪心结果全部保留。

\FloatBarrier\begingroup\small\setlength{\tabcolsep}{2.5pt}
\begin{longtable}[]{@{}
  >{\raggedright\arraybackslash}p{(\linewidth - 4\tabcolsep) * \real{0.3344}}
  >{\raggedleft\arraybackslash}p{(\linewidth - 4\tabcolsep) * \real{0.3292}}
  >{\raggedleft\arraybackslash}p{(\linewidth - 4\tabcolsep) * \real{0.3363}}@{}}
\caption{验证集阈值校准：设置与结果}\label{tab:s11-3}\tabularnewline
\toprule\noalign{}
\begin{minipage}[b]{\linewidth}\raggedright
校准后完整240例测试
\end{minipage} & \begin{minipage}[b]{\linewidth}\raggedleft
基座
\end{minipage} & \begin{minipage}[b]{\linewidth}\raggedleft
LoRA
\end{minipage} \\
\midrule\noalign{}
\endfirsthead
\toprule\noalign{}
\begin{minipage}[b]{\linewidth}\raggedright
校准后完整240例测试
\end{minipage} & \begin{minipage}[b]{\linewidth}\raggedleft
基座
\end{minipage} & \begin{minipage}[b]{\linewidth}\raggedleft
LoRA
\end{minipage} \\
\midrule\noalign{}
\endhead
\bottomrule\noalign{}
\endlastfoot
完成召回 & 3/39（7.7\%） & 5/39（12.8\%） \\
未完成误报 & 10/201（5.0\%） & 14/201（7.0\%） \\
常规负例误报 & 2/121（1.7\%） & 1/121（0.8\%） \\
边界80例误报 & 8/80（10.0\%） & 13/80（16.3\%） \\
\end{longtable}\endgroup

LoRA阈值为0.985936。验证集的5\%约束没有保证测试集也满足，且低误报设置损失了大部分完成召回。因此这轮权重尚不适合用于自动阶段切换。数据正例主要是成功末帧，可能包含执行阶段和机械臂姿态的关联；下一轮应补充外观相似但失败的末帧及反事实负例，而不是立即接入机器人。

基座和LoRA使用同一份该阶段提示词；这些数字不等同于旧分层系统原完成检查提示词的结果。仅含负例的边界子集不报告平衡准确率，避免给单类样本赋予误导性的数值。


\subsection{第二版检查器：近似完成数据与独立测试}\label{s11-detail-7}

\subsubsection{目标与验收}\label{s11-detail-8}

希望解决的具体问题是：物体看起来已经到达容器或盘子附近，但实际上仍被夹持、悬空或偏离有效区域时，检查器过早确认完成。

离线验收预先要求：未完成样例误报率不超过 5\%，完成召回至少 20\%，且较同输入形式的基座提高至少 15 个百分点。阈值由验证集选择，测试集不参与训练、选权重和阈值调整。未知或拒绝确认不计作完成召回。

离线通过后才运行配对闭环：固定既有计划和控制参数，对比基座与检查器，要求成功率不低于 28/30，同时减少错误切换。目前这一阶段未启动。

\subsubsection{数据与标签}\label{s11-detail-9}

\FloatBarrier\begingroup\small\setlength{\tabcolsep}{2.5pt}
\begin{longtable}[]{@{}
  >{\raggedright\arraybackslash}p{(\linewidth - 6\tabcolsep) * \real{0.2397}}
  >{\raggedleft\arraybackslash}p{(\linewidth - 6\tabcolsep) * \real{0.2397}}
  >{\raggedleft\arraybackslash}p{(\linewidth - 6\tabcolsep) * \real{0.2397}}
  >{\raggedright\arraybackslash}p{(\linewidth - 6\tabcolsep) * \real{0.2810}}@{}}
\caption{数据与标签：设置与结果}\label{tab:s11-4}\tabularnewline
\toprule\noalign{}
\begin{minipage}[b]{\linewidth}\raggedright
划分
\end{minipage} & \begin{minipage}[b]{\linewidth}\raggedleft
样例
\end{minipage} & \begin{minipage}[b]{\linewidth}\raggedleft
场景组
\end{minipage} & \begin{minipage}[b]{\linewidth}\raggedright
用途
\end{minipage} \\
\midrule\noalign{}
\endfirsthead
\toprule\noalign{}
\begin{minipage}[b]{\linewidth}\raggedright
划分
\end{minipage} & \begin{minipage}[b]{\linewidth}\raggedleft
样例
\end{minipage} & \begin{minipage}[b]{\linewidth}\raggedleft
场景组
\end{minipage} & \begin{minipage}[b]{\linewidth}\raggedright
用途
\end{minipage} \\
\midrule\noalign{}
\endhead
\bottomrule\noalign{}
\endlastfoot
训练 & 468 & 77 & 参数更新 \\
验证 & 288 & 54 & 选权重、阈值 \\
独立测试 & 284 & 53 & 一次性评价语言层模型 \\
\end{longtable}\endgroup

799 例来自仿真状态构造，241 例来自既有策略回合。合成场景使用 Long 初态 3--10；训练、验证和测试按初态及场景组隔离。既有 Long 初态 0--2 留作闭环评价。自然回合按来源划分，并检查图像对和场景组交叉。

合成样例对照使用相同的机器人姿态、相机和背景，避免模型只按机械臂是否回位判断。部分场景先放好另一个物体，覆盖第二个子目标。相机位置和视场角有小幅变化；当前数据仍局限于 LIBERO 的资产和任务，不能据此声称具备真实场景泛化能力。

完成标签要求同时满足任务关系、非机器人支撑接触、未被实际夹持。仅满足仿真的 In 几何关系不够，因为它可能包含悬空物体。临界偏心、接触边缘、目标严重遮挡和视觉相同而物理状态不同的样例标为 unknown。所有物理标签仅用于离线数据和评价，不作为运行时模型输入。

审核期间修正了只统计碰撞几何体造成的可见性计数问题；修正、重新生成后才开始训练。对外形相同的两只壶，避免使用图像中无法确定的内部编号。详细规则见 DATA\_REVIEW.md。

\subsubsection{训练过程}\label{s11-detail-10}

基座为本地已有 Qwen3-VL-4B，BF16。最初以完整 JSON 文本进行监督训练，验证召回不足；随后改为直接优化 yes/no/unknown 三分类决策，减少固定输出格式对损失的占用。

语言层版本对 q/v 做 rank 8、alpha 16、dropout 0.05 的 LoRA，约 295 万可训练参数，学习率 5e-5，梯度累积 3。两个开发阶段各 360 次更新，第二阶段从第一阶段权重继续训练。权重仅按验证集选择。

追加视觉层适配时，保留语言层 LoRA，并加入视觉注意力 qkv，总计约 374 万可训练参数，再训练 360 次更新。其验证阈值采用 2\% 误报上限作为安全余量，测试验收仍为 5\%。额外准备初态 11、12 的新确认数据，以免把已评分的原测试集重复当作独立测试；该阶段视觉版本未通过验证，确认集保持未评分。

概率接近 1 时，浮点数会丢失排序差异，因此阈值采用 yes 相对 no/unknown 的 log-odds，使用双精度计算。阈值不是置信度百分比。

\subsubsection{独立测试结果}\label{s11-detail-11}

语言层版本由新进程重新加载权重后测试。两者分别使用验证集确定的阈值，输入形式相同。

\FloatBarrier\begingroup\small\setlength{\tabcolsep}{2.5pt}
\begin{longtable}[]{@{}
  >{\raggedright\arraybackslash}p{(\linewidth - 8\tabcolsep) * \real{0.1842}}
  >{\raggedleft\arraybackslash}p{(\linewidth - 8\tabcolsep) * \real{0.2192}}
  >{\raggedleft\arraybackslash}p{(\linewidth - 8\tabcolsep) * \real{0.2283}}
  >{\raggedleft\arraybackslash}p{(\linewidth - 8\tabcolsep) * \real{0.1842}}
  >{\raggedright\arraybackslash}p{(\linewidth - 8\tabcolsep) * \real{0.1842}}@{}}
\caption{独立测试结果：设置与结果}\label{tab:s11-5}\tabularnewline
\toprule\noalign{}
\begin{minipage}[b]{\linewidth}\raggedright
模型
\end{minipage} & \begin{minipage}[b]{\linewidth}\raggedleft
完成召回
\end{minipage} & \begin{minipage}[b]{\linewidth}\raggedleft
未完成误报率
\end{minipage} & \begin{minipage}[b]{\linewidth}\raggedleft
确认数量
\end{minipage} & \begin{minipage}[b]{\linewidth}\raggedright
结论
\end{minipage} \\
\midrule\noalign{}
\endfirsthead
\toprule\noalign{}
\begin{minipage}[b]{\linewidth}\raggedright
模型
\end{minipage} & \begin{minipage}[b]{\linewidth}\raggedleft
完成召回
\end{minipage} & \begin{minipage}[b]{\linewidth}\raggedleft
未完成误报率
\end{minipage} & \begin{minipage}[b]{\linewidth}\raggedleft
确认数量
\end{minipage} & \begin{minipage}[b]{\linewidth}\raggedright
结论
\end{minipage} \\
\midrule\noalign{}
\endhead
\bottomrule\noalign{}
\endlastfoot
基座 & 3/105，2.9\% & 1/179，0.6\% & 4/284 & 召回不足 \\
语言层 LoRA & 28/105，26.7\% & 12/179，6.7\% & 40/284 & 误报超标 \\
\end{longtable}\endgroup

召回增加 23.8 个百分点，但未通过全部验收条件。基座更保守，因此不能仅比较召回就宣称整体改善。样例存在组内相关性，上表是当前固定测试集上的结果，不作统计显著性或真实部署安全保证。

在原验证集上单纯收紧语言层阈值，误报率约束降至 2\% 时，召回也降至 14.8\%，不足以解决问题。加入视觉适配后，选中第 120 次更新；重新加载的验证召回为 14.8\%，误报率 1.7\%，相对基座增益 13.0 个百分点，仍未达到开发筛选要求。视觉版本没有独立测试成绩。


\subsection{检查器接口与文件}\label{s11-detail-12}

\subsubsection{接口与代码目录}\label{s11-detail-13}

\noindent 运行与目录示例见\hyperlink{s11-command-2}{附录D的对应条目}。\par

相关数据代码在 \path|checker_v2/scripts/generate.py| 和 audit\_data.py；相关训练代码在 train\_classification.py、train\_vision.py；相关评价代码在 protocol.py、final\_evaluate.py。完整原始数据保存在本地，源码包不包含图像、仿真状态和基座权重。

新推理接口已保留为显式开关，默认系统仍使用原检查器。规划继续使用基座；适配器只用于适用范围内的搬运完成检查。开关炉、门操作等未覆盖类别退回原实现。接口经过单元测试，但没有把它写成已经完成的闭环验证。


\stagefigure{figures/ch04/s11_validation_false_positives.png}{检查器的误确认样例。最终版需逐项解释物理状态与视觉可辨性。}{fig:s11-retained}

\paragraph{阶段结论}V2整理1040例数据，完成召回从2.9\%提高到26.7\%；误报率6.7\%，超过预设5\%上限。视觉层适配没有通过验证集筛选。 误报率是未完成样例被误判完成的比例；召回率是已完成样例被正确确认的比例。离线结果不能直接替代闭环验证。

\stagecode{s11}

\FloatBarrier

\section{检查器的闭环作用与输入对照}

\label{stage:s12}

在真实执行回合中比较视觉检查、冻结图像、固定时刻切换和始终保持第一条子目标，并补充腕部视角与任务约束透传。

\subsection{冻结V2检查器的配对闭环}\label{s12-detail-1}

\subsubsection{为什么继续跑闭环}\label{s12-detail-2}

此前独立离线测试中，V2 将完成召回从 2.9\% 提高到 26.7\%，误报由 0.6\% 增至 6.7\%，未达到预设 5\% 上限。该阶段保留这一结论，但将它与探索性仿真分开：实际闭环还有夹爪门控、连续两次确认和低层策略行为，离线误报不能直接换算为任务失败。

离线标签要求到位、支撑接触且已释放，闭环则保留 LIBERO 官方成功检查。两套标准不同。该阶段不改离线标签、不重新选阈值，也不把原测试重新评分后当成新测试。

\subsubsection{固定的对照条件}\label{s12-detail-3}

\begin{itemize}
\tightlist
\item
  低层均为同一个官方 \(\pi\)0.5 LIBERO 检查点，10 步去噪，每次执行动作块的前 5 步。
\item
  固定 v3 已记录的英文子目标及中文整句指令；每回合固定初态、环境种子和动作采样噪声，不重新生成计划。
\item
  第三人称 512 像素开始帧与当前帧；夹爪张开时检查，连续两次确认才切换。
\item
  每子目标预算 300 步，整回合最多 520 步；末子目标由官方成功或步数上限结束。超预算推进与视觉确认分别计数。
\item
  V2 使用语言层已冻结适配器，log-odds 阈值 3.2480709186546264。门、炉开关和独立拿起等未覆盖步骤保留基座回退。
\item
  两组按任务与初态交替先后运行。仿真真值仅保存作评价，不参与检查或控制。
\end{itemize}

30 对计划和种子一致，首次切换前动作最大差为 0；其中 20 对全程动作完全相同。该阶段不是实体机械臂实验。

\subsubsection{闭环结果}\label{s12-detail-4}

\FloatBarrier\begingroup\small\setlength{\tabcolsep}{2.5pt}
\begin{longtable}[]{@{}
  >{\raggedright\arraybackslash}p{(\linewidth - 4\tabcolsep) * \real{0.3926}}
  >{\raggedleft\arraybackslash}p{(\linewidth - 4\tabcolsep) * \real{0.3037}}
  >{\raggedleft\arraybackslash}p{(\linewidth - 4\tabcolsep) * \real{0.3037}}@{}}
\caption{闭环结果：设置与结果}\label{tab:s12-1}\tabularnewline
\toprule\noalign{}
\begin{minipage}[b]{\linewidth}\raggedright
指标
\end{minipage} & \begin{minipage}[b]{\linewidth}\raggedleft
基座检查器
\end{minipage} & \begin{minipage}[b]{\linewidth}\raggedleft
V2 分类检查器
\end{minipage} \\
\midrule\noalign{}
\endfirsthead
\toprule\noalign{}
\begin{minipage}[b]{\linewidth}\raggedright
指标
\end{minipage} & \begin{minipage}[b]{\linewidth}\raggedleft
基座检查器
\end{minipage} & \begin{minipage}[b]{\linewidth}\raggedleft
V2 分类检查器
\end{minipage} \\
\midrule\noalign{}
\endhead
\bottomrule\noalign{}
\endlastfoot
官方成功 & 28/30 & 28/30 \\
超时回合 & 2 & 2 \\
视觉确认推进 & 19 & 12 \\
可核查的正确确认 & 17 & 10 \\
错误确认切换 & 1 & 1 \\
无对应最终谓词的确认 & 1 & 1 \\
超预算推进 & 2 & 1 \\
达成但未确认的回合/阶段 & 8 & 15 \\
成功但未发生子目标切换 & 9 & 17 \\
全部检查次数 & 135 & 145 \\
\end{longtable}\endgroup

达成但未确认包含最终任务已成功的情况，因此不等于任务失败。独立拿起书本没有对应的最终任务谓词，不能用``书已进收纳盒''替代拿起动作的真值；这类确认不进入错误切换计数。两只外观相同的壶，首阶段按任意一只已上炉评价。

仅统计正确确认的平均延迟分别为 2.13 / 2.14 秒仿真时间，样本数为 17 / 10。两组都正确确认的配对回合共 9 个，其平均延迟为 2.03 / 2.26 秒，V2 减基座为 0.22 秒。未确认项单列，避免只统计快速成功的案例。

配对成功结果：仅基座成功 0 回合，仅 V2 成功 0 回合；精确 McNemar 双侧 p=1。每任务仅 3 个初态，不据此作广泛泛化结论。

\FloatBarrier\begingroup\small\setlength{\tabcolsep}{2.5pt}
\begin{longtable}[]{@{}
  >{\raggedright\arraybackslash}p{(\linewidth - 4\tabcolsep) * \real{0.3677}}
  >{\raggedleft\arraybackslash}p{(\linewidth - 4\tabcolsep) * \real{0.3162}}
  >{\raggedleft\arraybackslash}p{(\linewidth - 4\tabcolsep) * \real{0.3162}}@{}}
\caption{闭环结果：设置与结果}\label{tab:s12-2}\tabularnewline
\toprule\noalign{}
\begin{minipage}[b]{\linewidth}\raggedright
任务
\end{minipage} & \begin{minipage}[b]{\linewidth}\raggedleft
基座
\end{minipage} & \begin{minipage}[b]{\linewidth}\raggedleft
V2
\end{minipage} \\
\midrule\noalign{}
\endfirsthead
\toprule\noalign{}
\begin{minipage}[b]{\linewidth}\raggedright
任务
\end{minipage} & \begin{minipage}[b]{\linewidth}\raggedleft
基座
\end{minipage} & \begin{minipage}[b]{\linewidth}\raggedleft
V2
\end{minipage} \\
\midrule\noalign{}
\endhead
\bottomrule\noalign{}
\endlastfoot
t0 双罐入篮 & 3/3 & 3/3 \\
t1 芝士黄油 & 3/3 & 3/3 \\
t2 开炉放壶 & 2/3 & 2/3 \\
t3 碗入抽屉 & 3/3 & 3/3 \\
t4 双杯放盘 & 3/3 & 3/3 \\
t5 书入收纳 & 3/3 & 3/3 \\
t6 杯与布丁 & 3/3 & 3/3 \\
t7 汤罐芝士 & 3/3 & 3/3 \\
t8 双壶上炉 & 2/3 & 2/3 \\
t9 杯入微波炉 & 3/3 & 3/3 \\
\end{longtable}\endgroup

\stagefigure{figures/ch04/s12_exploratory_success_bc9ca426.png}{各任务配对成功回合数}{fig:s12-detail-1}

\subsubsection{这轮验证了什么}\label{s12-detail-5}

第一，闭环可以在离线门槛未通过时作为探索继续运行，但两种结论应分别保留。该阶段既没有证明单次误报会直接造成闭环失败，也没有证明离线召回提高就能改善任务成功。

第二，一部分回合中，\(\pi\)0.5 在首个子目标尚未被确认时就完成了整项任务。这说明当前基准上的低层行为可能覆盖比单条子目标更长的动作序列。单看任务成功率不足以证明高层完成检查有效，因此该阶段补充了确认次数、未确认情况和全程动作相同比例。

第三，V2 的主要可测变化是检查方式和运行耗时；是否值得部署仍取决于对具体错误的改善，而不是权重已经训练出来。未覆盖操作的基座回退也必须单列，不能把它们当成 LoRA 已解决的类别。

\subsubsection{误报标签审核与后续取舍}\label{s12-detail-6}

原离线测试 12 个误报中，7 个是近接触、2 个明显偏移、2 个边缘外侧、1 个夹持未放；其中 9 个标签为 unknown，5 个仿真任务关系已成立但未满足更严格的支撑/释放要求。这是评分后的语义审核，不用于修改原阈值或擦除失败结果。

当前不再盲目扩大到相似的数千张静态图。若继续训练，应先针对确实影响闭环的错误补充自然过程和可辨识的困难对照；腕部相机先检查可见性再决定加入。完成检查器这一阶段按该阶段结果收尾，后续优先补同一个物理场景中的导航与抓取，沿用已有检查与记录接口。


\subsection{视觉信息干预与阶段切换归因}\label{s12-detail-7}

\subsubsection{问题与已有基础}\label{s12-detail-8}

原\(\pi\)0.5在LIBERO完成117/120次闭环；已有中文Long对照中，Qwen分层28/30、仅翻译29/30。已有检查器V2在探索性配对闭环与基座均28/30，单次检查更快，但没有改善成功率。这些结果说明系统运行正常，尚不能证明视觉阶段检查带来了额外成功。

该阶段的问题是：在不改动低层策略与计划的情况下，拿走当前视觉或取消阶段推进，任务是否仍能完成？若能，下一步需要设计真正需要高层介入的任务，而不是只在相同成功率上继续优化检查器。

既有数字只作背景，不与该阶段回合混算。当前30个初态是已经研究过的Long任务0--9、初态0--2，属于机制诊断，不是新的泛化评测。

\subsubsection{配对协议与五组干预}\label{s12-detail-9}

\FloatBarrier\begingroup\small\setlength{\tabcolsep}{2.5pt}
\begin{longtable}[]{@{}
  >{\raggedright\arraybackslash}p{(\linewidth - 4\tabcolsep) * \real{0.2281}}
  >{\raggedright\arraybackslash}p{(\linewidth - 4\tabcolsep) * \real{0.3135}}
  >{\raggedright\arraybackslash}p{(\linewidth - 4\tabcolsep) * \real{0.4584}}@{}}
\caption{配对协议与五组干预：设置与结果}\label{tab:s12-3}\tabularnewline
\toprule\noalign{}
\begin{minipage}[b]{\linewidth}\raggedright
组别
\end{minipage} & \begin{minipage}[b]{\linewidth}\raggedright
检查器可见内容
\end{minipage} & \begin{minipage}[b]{\linewidth}\raggedright
切换规则
\end{minipage} \\
\midrule\noalign{}
\endfirsthead
\toprule\noalign{}
\begin{minipage}[b]{\linewidth}\raggedright
组别
\end{minipage} & \begin{minipage}[b]{\linewidth}\raggedright
检查器可见内容
\end{minipage} & \begin{minipage}[b]{\linewidth}\raggedright
切换规则
\end{minipage} \\
\midrule\noalign{}
\endhead
\bottomrule\noalign{}
\endlastfoot
正常视觉 & 子目标开始帧＋真实当前帧 & 基座Qwen判断、释放门控、连续两票 \\
冻结旧图 & 开始帧＋同一开始帧 & 相同判断与门控；仅替换检查图像 \\
预设时刻 & 不调用Qwen & 150步后给完成票，仍需相同门控和两票 \\
V2检查方案 & 开始帧＋真实当前帧 & 原LoRA＋提示词＋分类阈值；不支持的操作回退基座 \\
保持首条指令 & 不调用Qwen & 不推进阶段，关闭300步阶段预算；直到官方成功或520步 \\
\end{longtable}\endgroup

前四组都保留300步阶段预算，超时推进与判断完成分别计数。预设150步是统一事前参数，不按任务调节；实际切换仍等待释放门控和两次查询，因此不一定在150步发生。保持首条指令是明确不同的机制控制，不能称为把完整长任务直接交给策略的flat基线。

所有方法都使用同一段冻结的历史Qwen计划，该阶段不再调用规划器生成新计划。因此不能把结果解释为中文理解、翻译或分解质量对照。中文整句保留在记录中，低层实际接收的是相同英文子目标。

策略使用官方\(\pi\)0.5检查点，每次预测10步、执行前5步再观察；显式输入相同种子的噪声。Qwen3-VL-4B和V2适配器在全部方法运行期间共驻，基座生成时关闭适配器。V2对照是提示词、LoRA和评分阈值整个方案，不单独归因为LoRA。

\stagefigure{figures/ch04/s12_image_intervention_0989cbe3.png}{冻结干预只改变检查输入。低层仍读取真实当前画面。}{fig:s12-detail-2}

相关代码在 \path|monitor_causality/scripts/interventions.py|、\path|run_study.py| 和 \path|runner_snapshot.py|。旧入口和历史结果没有改动。

\subsubsection{官方结果与任务差异}\label{s12-detail-10}

\FloatBarrier\begingroup\small\setlength{\tabcolsep}{2.5pt}
\begin{longtable}[]{@{}
  >{\raggedright\arraybackslash}p{(\linewidth - 8\tabcolsep) * \real{0.2036}}
  >{\raggedright\arraybackslash}p{(\linewidth - 8\tabcolsep) * \real{0.1991}}
  >{\raggedright\arraybackslash}p{(\linewidth - 8\tabcolsep) * \real{0.1991}}
  >{\raggedright\arraybackslash}p{(\linewidth - 8\tabcolsep) * \real{0.1991}}
  >{\raggedright\arraybackslash}p{(\linewidth - 8\tabcolsep) * \real{0.1991}}@{}}
\caption{官方结果与任务差异：设置与结果}\label{tab:s12-4}\tabularnewline
\toprule\noalign{}
\begin{minipage}[b]{\linewidth}\raggedright
方法
\end{minipage} & \begin{minipage}[b]{\linewidth}\raggedright
官方成功
\end{minipage} & \begin{minipage}[b]{\linewidth}\raggedright
成功且未切换
\end{minipage} & \begin{minipage}[b]{\linewidth}\raggedright
实际Qwen检查调用
\end{minipage} & \begin{minipage}[b]{\linewidth}\raggedright
阶段切换总数
\end{minipage} \\
\midrule\noalign{}
\endfirsthead
\toprule\noalign{}
\begin{minipage}[b]{\linewidth}\raggedright
方法
\end{minipage} & \begin{minipage}[b]{\linewidth}\raggedright
官方成功
\end{minipage} & \begin{minipage}[b]{\linewidth}\raggedright
成功且未切换
\end{minipage} & \begin{minipage}[b]{\linewidth}\raggedright
实际Qwen检查调用
\end{minipage} & \begin{minipage}[b]{\linewidth}\raggedright
阶段切换总数
\end{minipage} \\
\midrule\noalign{}
\endhead
\bottomrule\noalign{}
\endlastfoot
正常视觉 & 28/30 & 9 & 132 & 21 \\
冻结旧图 & 28/30 & 25 & 157 & 5 \\
预设时刻 & 28/30 & 10 & 0 & 20 \\
V2检查方案 & 28/30 & 17 & 139 & 13 \\
保持首条指令 & 28/30 & 28 & 0 & 0 \\
\end{longtable}\endgroup

\FloatBarrier\begingroup\small\setlength{\tabcolsep}{2.5pt}
\begin{longtable}[]{@{}
  >{\raggedright\arraybackslash}p{(\linewidth - 10\tabcolsep) * \real{0.1667}}
  >{\raggedright\arraybackslash}p{(\linewidth - 10\tabcolsep) * \real{0.1667}}
  >{\raggedright\arraybackslash}p{(\linewidth - 10\tabcolsep) * \real{0.1667}}
  >{\raggedright\arraybackslash}p{(\linewidth - 10\tabcolsep) * \real{0.1667}}
  >{\raggedright\arraybackslash}p{(\linewidth - 10\tabcolsep) * \real{0.1667}}
  >{\raggedright\arraybackslash}p{(\linewidth - 10\tabcolsep) * \real{0.1667}}@{}}
\caption{官方结果与任务差异：设置与结果}\label{tab:s12-5}\tabularnewline
\toprule\noalign{}
\begin{minipage}[b]{\linewidth}\raggedright
任务
\end{minipage} & \begin{minipage}[b]{\linewidth}\raggedright
正常视觉
\end{minipage} & \begin{minipage}[b]{\linewidth}\raggedright
冻结旧图
\end{minipage} & \begin{minipage}[b]{\linewidth}\raggedright
预设时刻
\end{minipage} & \begin{minipage}[b]{\linewidth}\raggedright
V2检查
\end{minipage} & \begin{minipage}[b]{\linewidth}\raggedright
保持首条
\end{minipage} \\
\midrule\noalign{}
\endfirsthead
\toprule\noalign{}
\begin{minipage}[b]{\linewidth}\raggedright
任务
\end{minipage} & \begin{minipage}[b]{\linewidth}\raggedright
正常视觉
\end{minipage} & \begin{minipage}[b]{\linewidth}\raggedright
冻结旧图
\end{minipage} & \begin{minipage}[b]{\linewidth}\raggedright
预设时刻
\end{minipage} & \begin{minipage}[b]{\linewidth}\raggedright
V2检查
\end{minipage} & \begin{minipage}[b]{\linewidth}\raggedright
保持首条
\end{minipage} \\
\midrule\noalign{}
\endhead
\bottomrule\noalign{}
\endlastfoot
0 & 3/3 & 3/3 & 3/3 & 3/3 & 3/3 \\
1 & 3/3 & 3/3 & 3/3 & 3/3 & 3/3 \\
2 & 2/3 & 2/3 & 2/3 & 2/3 & 2/3 \\
3 & 3/3 & 3/3 & 3/3 & 3/3 & 3/3 \\
4 & 3/3 & 3/3 & 3/3 & 3/3 & 3/3 \\
5 & 3/3 & 3/3 & 3/3 & 3/3 & 3/3 \\
6 & 3/3 & 3/3 & 3/3 & 3/3 & 3/3 \\
7 & 3/3 & 3/3 & 3/3 & 3/3 & 3/3 \\
8 & 2/3 & 2/3 & 2/3 & 2/3 & 2/3 \\
9 & 3/3 & 3/3 & 3/3 & 3/3 & 3/3 \\
\end{longtable}\endgroup

任务编号对应如下，完整英文子目标逐初态保留在PROTOCOL.json：

\FloatBarrier\begingroup\small\setlength{\tabcolsep}{2.5pt}
\begin{longtable}[]{@{}
  >{\raggedright\arraybackslash}p{(\linewidth - 2\tabcolsep) * \real{0.3156}}
  >{\raggedright\arraybackslash}p{(\linewidth - 2\tabcolsep) * \real{0.6844}}@{}}
\caption{官方结果与任务差异：设置与结果}\label{tab:s12-6}\tabularnewline
\toprule\noalign{}
\begin{minipage}[b]{\linewidth}\raggedright
任务
\end{minipage} & \begin{minipage}[b]{\linewidth}\raggedright
长任务指令
\end{minipage} \\
\midrule\noalign{}
\endfirsthead
\toprule\noalign{}
\begin{minipage}[b]{\linewidth}\raggedright
任务
\end{minipage} & \begin{minipage}[b]{\linewidth}\raggedright
长任务指令
\end{minipage} \\
\midrule\noalign{}
\endhead
\bottomrule\noalign{}
\endlastfoot
0 & 把字母汤罐头和番茄沙司罐头都放进篮子里 \\
1 & 把奶油奶酪盒和黄油都放进篮子里 \\
2 & 打开炉子，然后把摩卡壶放到炉子上 \\
3 & 把黑色的碗放进柜子最下面的抽屉，然后关上抽屉 \\
4 & 把白色马克杯放到左边的盘子上，再把黄白相间的马克杯放到右边的盘子上 \\
5 & 拿起那本书，放进桌面收纳盒靠后的格子里 \\
6 & 把白色马克杯放到盘子上，再把巧克力布丁放到盘子右边 \\
7 & 把字母汤罐头和奶油奶酪盒都放进篮子里 \\
8 & 把两个摩卡壶都放到炉子上 \\
9 & 把黄白相间的马克杯放进微波炉，然后关上微波炉门 \\
\end{longtable}\endgroup

同一初态中，以正常视觉为左侧参照，另一方法增加/丢失的成功如下：

\FloatBarrier\begingroup\small\setlength{\tabcolsep}{2.5pt}
\begin{longtable}[]{@{}
  >{\raggedright\arraybackslash}p{(\linewidth - 10\tabcolsep) * \real{0.1769}}
  >{\raggedright\arraybackslash}p{(\linewidth - 10\tabcolsep) * \real{0.1646}}
  >{\raggedright\arraybackslash}p{(\linewidth - 10\tabcolsep) * \real{0.1646}}
  >{\raggedright\arraybackslash}p{(\linewidth - 10\tabcolsep) * \real{0.1646}}
  >{\raggedright\arraybackslash}p{(\linewidth - 10\tabcolsep) * \real{0.1646}}
  >{\raggedright\arraybackslash}p{(\linewidth - 10\tabcolsep) * \real{0.1646}}@{}}
\caption{官方结果与任务差异：设置与结果}\label{tab:s12-7}\tabularnewline
\toprule\noalign{}
\begin{minipage}[b]{\linewidth}\raggedright
对照方法
\end{minipage} & \begin{minipage}[b]{\linewidth}\raggedright
配对数
\end{minipage} & \begin{minipage}[b]{\linewidth}\raggedright
增加成功
\end{minipage} & \begin{minipage}[b]{\linewidth}\raggedright
丢失成功
\end{minipage} & \begin{minipage}[b]{\linewidth}\raggedright
名义精确McNemar p
\end{minipage} & \begin{minipage}[b]{\linewidth}\raggedright
共同成功时平均步数差
\end{minipage} \\
\midrule\noalign{}
\endfirsthead
\toprule\noalign{}
\begin{minipage}[b]{\linewidth}\raggedright
对照方法
\end{minipage} & \begin{minipage}[b]{\linewidth}\raggedright
配对数
\end{minipage} & \begin{minipage}[b]{\linewidth}\raggedright
增加成功
\end{minipage} & \begin{minipage}[b]{\linewidth}\raggedright
丢失成功
\end{minipage} & \begin{minipage}[b]{\linewidth}\raggedright
名义精确McNemar p
\end{minipage} & \begin{minipage}[b]{\linewidth}\raggedright
共同成功时平均步数差
\end{minipage} \\
\midrule\noalign{}
\endhead
\bottomrule\noalign{}
\endlastfoot
冻结旧图 & 30 & 0 & 0 & 1.000 & +0.1 \\
预设时刻 & 30 & 0 & 0 & 1.000 & -3.2 \\
V2检查方案 & 30 & 0 & 0 & 1.000 & -1.2 \\
保持首条指令 & 30 & 0 & 0 & 1.000 & -0.0 \\
\end{longtable}\endgroup

p值仅描述配对不一致项；每任务只有3个初态，任务内相关，且同时比较多个方法，不能据此宣称一般等效或显著优越。步数差为对照方法减正常视觉，只比较双方都成功的回合，不是墙钟耗时。相同成功数也可能对应不同回合或不同动作轨迹。

不同结局的配对条件（步数只描述该回合，不用失败提前结束冒充速度提升）：

该阶段30个配对条件中，五种方法的成功/失败结局逐项一致。

\subsubsection{判断、切换与低层行为}\label{s12-detail-11}

\paragraph{不把预算推进当作视觉确认}\label{s12-detail-12}

错误切换按切换当时的官方中间任务关系判断，预算到期的错误推进另计。任务5的第一步是独立拿书，最终任务谓词不能直接评估``已拿起''，因此标为不可评分，不强行算对或错。

\FloatBarrier\begingroup\small\setlength{\tabcolsep}{2.5pt}
\begin{longtable}[]{@{}
  >{\raggedright\arraybackslash}p{(\linewidth - 10\tabcolsep) * \real{0.1698}}
  >{\raggedright\arraybackslash}p{(\linewidth - 10\tabcolsep) * \real{0.1660}}
  >{\raggedright\arraybackslash}p{(\linewidth - 10\tabcolsep) * \real{0.1660}}
  >{\raggedright\arraybackslash}p{(\linewidth - 10\tabcolsep) * \real{0.1660}}
  >{\raggedright\arraybackslash}p{(\linewidth - 10\tabcolsep) * \real{0.1660}}
  >{\raggedright\arraybackslash}p{(\linewidth - 10\tabcolsep) * \real{0.1660}}@{}}
\caption{不把预算推进当作视觉确认：设置与结果}\label{tab:s12-8}\tabularnewline
\toprule\noalign{}
\begin{minipage}[b]{\linewidth}\raggedright
方法
\end{minipage} & \begin{minipage}[b]{\linewidth}\raggedright
完成票触发的错误切换
\end{minipage} & \begin{minipage}[b]{\linewidth}\raggedright
预算推进时未完成
\end{minipage} & \begin{minipage}[b]{\linewidth}\raggedright
不可评分切换
\end{minipage} & \begin{minipage}[b]{\linewidth}\raggedright
正确切换平均延迟（步）
\end{minipage} & \begin{minipage}[b]{\linewidth}\raggedright
平均模型检查耗时（s）
\end{minipage} \\
\midrule\noalign{}
\endfirsthead
\toprule\noalign{}
\begin{minipage}[b]{\linewidth}\raggedright
方法
\end{minipage} & \begin{minipage}[b]{\linewidth}\raggedright
完成票触发的错误切换
\end{minipage} & \begin{minipage}[b]{\linewidth}\raggedright
预算推进时未完成
\end{minipage} & \begin{minipage}[b]{\linewidth}\raggedright
不可评分切换
\end{minipage} & \begin{minipage}[b]{\linewidth}\raggedright
正确切换平均延迟（步）
\end{minipage} & \begin{minipage}[b]{\linewidth}\raggedright
平均模型检查耗时（s）
\end{minipage} \\
\midrule\noalign{}
\endhead
\bottomrule\noalign{}
\endlastfoot
正常视觉 & 0 & 2 & 1 & 45.7 & 0.521 \\
冻结旧图 & 0 & 2 & 0 & - & 0.364 \\
预设时刻 & 0 & 2 & 0 & 52.6 & - \\
V2检查方案 & 0 & 2 & 1 & 33.7 & 0.227 \\
保持首条指令 & 0 & 0 & 0 & - & - \\
\end{longtable}\endgroup

正确切换延迟从切换前连续保持为真的最近起点计算。曾经达成又丢失的状态不会被算作正确切换；平均值只覆盖实际正确切换，未确认状态单独记录，不能用较少确认造成的低延迟``证明更好''。计时组的完成票是规则决定，不是视觉判断。

这里的两票是连续两次通过门控的检查，不是相邻视频帧。例如t2初态1在第40步被Qwen判断为已开炉，而官方关系仍为假；下一次符合门控的检查已到第120步，此时炉子确实开启，第二票才触发推进。中间80步没有新检查，旧规则也没有票据过期时间。该阶段保留该规则，因此不能把两票解释为两份独立判断。

\stagefigure{figures/ch04/s12_decisions_cost_76ce9d0d.png}{切换错误与实际Qwen调用成本。}{fig:s12-detail-3}

\paragraph{检查无贡献与策略不看视觉是两件事}\label{s12-detail-13}

冻结组只拿走Qwen的当前检查画面，\(\pi\)0.5仍以当前RGB、腕部图和机器人状态闭环预测动作。因而冻结组成功不能解释为低层不用视觉，也不表示VLA是动作回放。

保持首条指令组的28/30成功说明：至少在这些初态中，推进到第二条指令不是达到官方长任务目标的必要条件。可能原因包括低层对已知场景的整任务先验、子目标指令本身覆盖更长行为等；该阶段不能仅凭成功率确定模型内部原因。

相反，某些回合如需切换才能成功，也应逐项定位，不能被总体计数掩盖。全量逐回合结果、动作与检查事件保留，便于分析不同任务。

再按每个动作维度逐元素比较整段策略输出，下面统计与正常视觉组完全一致的回合：

\FloatBarrier\begingroup\small\setlength{\tabcolsep}{2.5pt}
\begin{longtable}[]{@{}
  >{\raggedright\arraybackslash}p{(\linewidth - 4\tabcolsep) * \real{0.3393}}
  >{\raggedright\arraybackslash}p{(\linewidth - 4\tabcolsep) * \real{0.3303}}
  >{\raggedright\arraybackslash}p{(\linewidth - 4\tabcolsep) * \real{0.3303}}@{}}
\caption{检查无贡献与策略不看视觉是两件事：设置与结果}\label{tab:s12-9}\tabularnewline
\toprule\noalign{}
\begin{minipage}[b]{\linewidth}\raggedright
对照
\end{minipage} & \begin{minipage}[b]{\linewidth}\raggedright
整段执行动作完全相同
\end{minipage} & \begin{minipage}[b]{\linewidth}\raggedright
配对总数
\end{minipage} \\
\midrule\noalign{}
\endfirsthead
\toprule\noalign{}
\begin{minipage}[b]{\linewidth}\raggedright
对照
\end{minipage} & \begin{minipage}[b]{\linewidth}\raggedright
整段执行动作完全相同
\end{minipage} & \begin{minipage}[b]{\linewidth}\raggedright
配对总数
\end{minipage} \\
\midrule\noalign{}
\endhead
\bottomrule\noalign{}
\endlastfoot
冻结旧图 & 10 & 30 \\
预设时刻 & 11 & 30 \\
V2检查方案 & 19 & 30 \\
保持首条指令 & 9 & 30 \\
\end{longtable}\endgroup

这把``官方结果一样''与``实际行为一样''分开。动作不同而同样成功，说明切换影响了行为，但尚未增加这些条件下的成功；整段动作相同，则该阶段该配对没有产生执行层差异。

\subsubsection{实际决策时刻的数据}\label{s12-detail-14}

整理728次实际检查机会，按参考RGB、当前RGB及指令去重后为162对：官方关系为真86、为假67、不可评分9，同像素标签冲突0对。

这与原先从任意时刻或成功末帧抽样不同，覆盖了真实释放门控之后系统准备作判断的状态。但尚未完成可见性审核，不能自动把仿真关系为真解释成人眼可以看清，也不能把它当成严格支撑、无夹持的物理标签。

索引保留原始参考/当前图、实际送入检查器的图、方法、时间步、输出和标签来源。冻结或计时组没有看真实当前图，其预测不能直接当成正常视觉误判来挖训练样本。任务5保持不可评分；严格物理支撑标签为空，可判定性待审。

所有采集仍来自既有初态0--2，划为开发采集，未随机拆成新的``独立测试集''。该阶段没有用它选择阈值、训练权重或报告新的离线模型提升。

相关代码在 \path|collect_queries.py|，数据索引在 \path|decision_data/cases.json|。

\subsubsection{独立核验与限制}\label{s12-detail-15}

\FloatBarrier\begingroup\small\setlength{\tabcolsep}{2.5pt}
\begin{longtable}[]{@{}
  >{\raggedright\arraybackslash}p{(\linewidth - 2\tabcolsep) * \real{0.5823}}
  >{\raggedright\arraybackslash}p{(\linewidth - 2\tabcolsep) * \real{0.4177}}@{}}
\caption{独立核验与限制：设置与结果}\label{tab:s12-10}\tabularnewline
\toprule\noalign{}
\begin{minipage}[b]{\linewidth}\raggedright
核验项目
\end{minipage} & \begin{minipage}[b]{\linewidth}\raggedright
结果
\end{minipage} \\
\midrule\noalign{}
\endfirsthead
\toprule\noalign{}
\begin{minipage}[b]{\linewidth}\raggedright
核验项目
\end{minipage} & \begin{minipage}[b]{\linewidth}\raggedright
结果
\end{minipage} \\
\midrule\noalign{}
\endhead
\bottomrule\noalign{}
\endlastfoot
完整回合 & 150 \\
完整解码视频帧 & 39452 \\
策略输入/噪声/动作块查询 & 7912 \\
检查机会 & 728 \\
切换前跨方法配对 & 120 \\
切换前最大动作差 & 0.0 \\
冻结组实际替换的不同当前画面 & 157 \\
协议、共享源码和适配器 & 哈希未变 \\
CPU单元检查 & 10项通过 \\
\end{longtable}\endgroup

逐查询核对RGB、腕部图、状态和噪声。每个实际执行动作都与保存的模型动作块前5步一致；相同任务/初态在最早切换之前的输入和动作逐项相同。冻结组还核对模型实际处理的两张图像素哈希相同，计时/保持组实际模型调用为0。

仿真中间谓词仅用于日志与离线评估，不送给检查器或策略来决定阶段。最终成功停止沿用官方LIBERO规则，因此不笼统声称``任何终止决策都不用真值''。

视频为实际回合的20fps仿真回放，不包含模型计算时的墙钟等待。五组对照片在某组结束后保留其真实末帧并标注已停止，不继续生成动作。完整视频解码核对通过。

主要限制是样本小、任务与初态已被研究过、低层检查点熟悉该基准，以及释放门控会改变检查的状态分布。这里诊断的是现有固定分层方案的作用，不能推广为``VLM高层没有用''。


\subsection{闭环复核入口}\label{s12-detail-16}

\subsubsection{代码与复核}\label{s12-detail-17}

\noindent 运行与目录示例见\hyperlink{s12-command-3}{附录D的对应条目}。\par

相关执行代码在 \path|scripts/evaluate_loop.py|；相关指标代码在 \path|scripts/summarize_loop.py|；检查器运行封装在 \path|experiments/vla_libero/scripts/completion_checker_v2.py|。完整 60 个回合的结果、视频和模型输入保存在本地，交付目录附代表视频，源码包不包含基座和训练图像。

\begin{Verbatim}[breaklines=true,breakanywhere=true,fontsize=\footnotesize]
# 已存在结果会复用，用于中断续跑；不覆盖原回合。
experiments/vla_libero/scripts/python_shared experiments/vla_libero/checker_v2/scripts/evaluate_loop.py --exploratory
.venvs/openpi-libero/bin/python experiments/vla_libero/checker_v2/scripts/summarize_loop.py
.venvs/openpi-libero/bin/python -m unittest discover -s experiments/vla_libero/checker_v2/scripts -p 'test_*.py'
\end{Verbatim}


\stagefigure{figures/ch04/s12_outcomes.png}{固定低层条件后的视觉干预对照，用于检验阶段检查的实际作用。}{fig:s12-retained}

\paragraph{阶段结论}五组主要干预共150回合；正常视觉、冻结图像和保持第一条指令均为28/30。当前任务不能充分证明视觉检查器的额外收益。 需要区分低层策略自身覆盖整项任务和高层阶段切换的作用。两次连续确认也可能重复同一种误判。

\stagecode{s12}

\FloatBarrier
